\begin{resumo}[Abstract]
 \begin{otherlanguage*}{english}
    The growing adoption of the Go programming language in modern architectures highlights a persistent challenge in software security: the high rate of false positives generated by Static Application Security Testing (SAST) tools. This scenario stems from the inherently conservative nature of static analyzers, which, in an attempt not to omit real threats, issue excessive alerts when examining code, especially when processing the architectural abstractions and concurrency models inherent to the Go language. To mitigate alert fatigue within development teams, this work aims to develop an automated approach for the triage and filtering of false positives in SAST reports through the design of a neuro-symbolic architecture. This solution seeks to integrate the structural precision of the code analysis engine with the semantic interpretation capabilities of Large Language Models (LLMs), employed as secondary classifiers. The methodology is based on a quantitative empirical experiment that couples the Semgrep engine with frontier neural models. The process utilizes the extraction of adjacent context and its subsequent injection into expert \textit{prompts}, which are structured with theoretical security rules and language-specific peculiarities. As a result, the guided reasoning of the AI is expected to filter analytical noise, reducing the cognitive effort required for manual validation without compromising the detection of genuine vulnerabilities.

   \vspace{\onelineskip}
 
   \noindent 
   \textbf{Key-words}: static analysis. Go language. artificial intelligence. automated triage. neuro-symbolic architecture. software security.
 \end{otherlanguage*}
\end{resumo}